\documentclass{situationist-preprint}

\title{Death and the Shape of a Life}
\msid{EXI-19-a}
\msfamily{EXI-19}
\msdiscipline{Meaning, Death and the Absurd}
\mstype{Principal manuscript}
\mskeywords{meaning of death; shape of a life; finitude; narrative; Velleman; Scheffler; Heidegger; deprivation}

\begin{document}
\maketitle

\begin{abstract}
What is the meaning of death? Two answers are common. On the first, death has a meaning of its own: it is a punishment, a passage, a completion or a test. On the second, death is meaningless and only takes away. We argue for a third answer. Death has no meaning in itself, since every candidate meaning either depends on claims that cannot be established or belongs to death only through its relation to the life it ends. But death gives life a shape by limiting it. Following Velleman, we take the value of a life to depend in part on its shape, the order in which its goods and evils occur, and we argue that a life has a shape as a whole only if it is bounded. The result is compatible with EXI-02: the limit that gives a life a shape is not a benefit to the one who dies, since a later death would also have given her life a shape, and in the ordinary case a better one. It complements EXI-03, which argues that an endless life is not desirable.
\end{abstract}

\section{The problem}

The question \enquote{What is the meaning of death?} asks for more than what death is. It asks what death signifies or contributes: whether it has a point, and whether it bears on the point of the life it ends. Religious traditions have given death a meaning of its own, as the wage of sin, as a passage to another life, or as the moment of judgment. Philosophers have more often found its meaning in its relation to life. Heidegger argued that it is only in being towards death that existence can be grasped as a whole \citep[\S\S 46--53]{Heidegger1927}; PHE-07 argues that such an authentic relation to death is possible as an orientation but not as a stable achievement. \citet{Nussbaum1989} argued that many human virtues and values, such as courage and the love of particular persons, have their character only for beings whose time is limited, and \citet{Scheffler2013} that the value of many of our activities depends on the scarcity of time. \citet{Velleman1991} showed that the value of a life is not the sum of the values of its moments, since a life that improves from misfortune to success is better than one that declines from success to misfortune with the same momentary goods. We combine the last two points: death has no meaning in itself, but it gives life a shape by limiting it.

\section{Framework}

\begin{definition}[Meaning in itself]\label{d:itself}
An event has \emph{meaning in itself} if it has a significance or point that it possesses independently of its relations to the life in which it occurs: if it is the fulfilment of a purpose, the communication of something, or the occasion of a value that it does not owe to that life.
\end{definition}

\begin{definition}[Shape]\label{d:shape}
The \emph{shape} of a life is the temporal order of its goods and evils, together with the relations among them that this order makes possible: progress, decline, completion, redemption and the like.
\end{definition}

\begin{principle}[Shape-dependence]\label{p:shape}
The value of a life as a whole depends in part on its shape, and not only on the sum of the momentary values it contains \citep{Velleman1991}.
\end{principle}

\begin{assumption}[Bounds]\label{a:bounds}
A sequence has an order \emph{as a whole}, of the kind Definition~\ref{d:shape} describes, only if it has a first and a last member. A sequence without a last member has shapes in each of its finite segments but no overall trajectory: no last state relative to which its course counts as progress, decline or completion.
\end{assumption}

On the tenseless view of time defended in TSM-01, a life is an ordered series of events, each located at a time. Assumption~\ref{a:bounds} says what is needed for that series to have the kind of order that makes it, as a whole, a story of rise or fall. A story that never ends may have many episodes, each with a shape; it has no shape as a story.

\section{The argument}

\begin{proposition}\label{p:noitself}
Death has no meaning in itself.
\end{proposition}

\begin{proof}
Consider the candidates. That death is a punishment, a passage or a judgment requires that it fulfil the purpose of an agent who imposes it or that it lead to a further life; neither can be established, and REL-08 argues that no evidence available before death could confirm survival. Even if one of these were true, the meaning would belong to death in virtue of its relation to a further life or a further purpose, not in itself. That death is a completion is a claim about its relation to the life it ends, and it is often false: most deaths interrupt projects rather than complete them. The remaining candidates, that death is natural or that it makes room for others, describe facts about death that are true independently of any life, but they assign it no significance or point: that something is natural, or causally necessary for something else, does not make it the fulfilment of a purpose or the bearer of a value. So no candidate meaning belongs to death in the sense of Definition~\ref{d:itself}.
\end{proof}

\begin{theorem}\label{t:main}
Death gives a life a shape by limiting it.
\end{theorem}

\begin{proof}
A life is an ordered series of events (Assumption~\ref{a:bounds} and the tenseless view). It begins with the person's coming to exist. If it ends, it ends with her death; so death supplies the last member of the series. By Assumption~\ref{a:bounds}, a series without a last member has no shape as a whole. So it is in virtue of being limited by death that a human life has a shape as a whole, and by Principle~\ref{p:shape} this shape contributes to the value of the life. Death does not determine what the shape is: that depends on the goods and evils the life contains and their order. It determines that there is one.
\end{proof}

\begin{corollary}\label{c:comp}
Theorem~\ref{t:main} is consistent with the claim of EXI-02 that death is bad for the one who dies when it deprives her of goods.
\end{corollary}

\begin{proof}
By the deprivation account of EXI-02, the value of a death for the one who dies is the difference between the value of her actual life and the value of the life she would have had had she not died then. In the closest world in which she does not die at that time, she dies later, so her life is limited and has a shape there too. The shaping function of Theorem~\ref{t:main} is therefore present on both sides of the comparison, and does not by itself make her actual death better for her than its alternative. In the ordinary case, the later death gives a better shape, since it allows projects to be completed that the earlier death interrupts.
\end{proof}

The result fits the conclusion of EXI-03 that an endless life is not desirable, but does not depend on it. EXI-03 argues that an endless life would exhaust the categorical desires that give life its point; Theorem~\ref{t:main} adds that it would have no shape as a whole. The two are independent: a life could lack a shape as a whole and still, for a long time, be full of point.

\section{Objections and replies}

\begin{objection}[Other limits]\label{o:other}
Projects, relationships and stages of life have beginnings and ends, and give life all the shape it needs. Death adds nothing.
\end{objection}

\begin{reply}
They give shape to the episodes of a life, and an endless life could have shaped episodes. What they do not give is the shape of the life as a whole, which is what Principle~\ref{p:shape} concerns: whether the life, taken as one thing, was a story of progress or of decline. Only the last limit settles that.
\end{reply}

\begin{objection}[Death as good]\label{o:good}
If shape is valuable and death confers it, then death is good for the one who dies, contrary to EXI-02.
\end{objection}

\begin{reply}
Corollary~\ref{c:comp} answers this. That a life is limited is a condition of its having a shape; that it is limited \emph{here} rather than later is what the deprivation account evaluates, and a later limit confers a shape too. Death in general is a condition of the good of shape; a particular death is usually a loss.
\end{reply}

\begin{objection}[Ruined shapes]\label{o:ruin}
Many deaths do not shape lives but spoil them: a young person dies before anything has taken form.
\end{objection}

\begin{reply}
Theorem~\ref{t:main} says that death gives a life a shape, not a good one. An early death gives a life the shape of a beginning cut off, which is one of the worst shapes a life can have; this is part of why, by EXI-02 and Principle~\ref{p:shape}, early death is so bad.
\end{reply}

\section{Conclusion}

Death has no meaning in itself. It is not the fulfilment of a purpose or the bearer of a value of its own, and the meanings attributed to it either cannot be established or belong to it through its relation to the life it ends. But it limits that life, and by limiting it gives it a shape as a whole, which is part of what makes a life go well or badly. This does not make any particular death good for the one who dies.

\begin{thebibliography}{99}\small
\bibitem[Heidegger(1927)]{Heidegger1927} Heidegger, M. (1927). \emph{Sein und Zeit}. Niemeyer. Trans. J. Macquarrie and E. Robinson, \emph{Being and Time}, Harper \& Row, 1962.
\bibitem[Nussbaum(1989)]{Nussbaum1989} Nussbaum, M.~C. (1989). Mortal immortals: Lucretius on death and the voice of nature. \emph{Philosophy and Phenomenological Research}, 50(2), 303--351.
\bibitem[Scheffler(2013)]{Scheffler2013} Scheffler, S. (2013). \emph{Death and the Afterlife}. Oxford University Press.
\bibitem[Velleman(1991)]{Velleman1991} Velleman, J.~D. (1991). Well-being and time. \emph{Pacific Philosophical Quarterly}, 72(1), 48--77.
\end{thebibliography}

\end{document}
